% Constructive discrete-complete GPSS complementation and exact endpoint scope.

\subsection{Constructive complementation of GPSS decompositions}
\label{sec:p09-gpss-decomposition-complement}

Use the exact future-labelled decomposition class of
\cref{def:p09-gpss-decomposition}. Point and gap labels are supplied
with their finite \(L(U)\) formulas; they are evaluated in the
whole model and have no individual parameter other than their
current point. This section proves relative complementation for
bounded intervals and for future rays with an actual lower endpoint.
It also settles why the same future-labelled class cannot support
the broader prefix and whole-order complement claim
\cite[Proposition~2.4 and its outline, p.~171]{GabbayPnueliShelahStavi1980}.

\begin{hblemma}{thm:p09-gpss-finite-bounded-interval}{proved=annex:p9.gpss.finite.bounded.interval}{Bounded intervals in discrete complete orders are finite}
In a nonempty discrete complete strict linear order, every closed
interval \([a,b]\), with actual endpoints \(a\leq b\),
is finite. Consequently every nonempty subset of an open bounded
interval has a first and last member. An adjacent pair has an empty
open interval. No finiteness of the entire order is asserted.
\end{hblemma}

\begin{hbdefinition}{def:p09-gpss-bounded-events}{A finite partial event language on a bounded interval}
Fix a discrete complete model and actual endpoints \(x<y\).
The root terms are \(x,y\).
For a future temporal label \(A\), form
\[
 \operatorname{First}_A(s,t),\qquad
 \operatorname{Last}_A(s,t),\qquad
 \operatorname{Test}_A(s).
\]
First and last return the first or last point satisfying \(A\)
strictly between the values of their two bounds. They are undefined
if a bound is undefined, the bounds are not strictly ordered, or the
occurrence set is empty. Tests are applied only to nonroot terms:
\(\operatorname{Test}_A(s)\) returns the value of \(s\) if
that value exists and satisfies \(A\), and is undefined otherwise.
A test is a first-occurrence lookup on the singleton consisting of
the already selected point.

Every term is a finite expression in earlier terms. Write
\(\operatorname{Def}(t)\) for its Boolean definedness status.
The undefined symbol is not a domain point and is never compared
with points by \(<\). Every defined nonroot term lies in
\((x,y)\). These are metalinguistic event terms, not additional
temporal connectives or predicates in the output formula. For example,
the source's last \(q\) before the first \(p\) in \((x,y)\)
is \(\operatorname{Last}_q(x,\operatorname{First}_p(x,y))\).
\end{hbdefinition}

\begin{hbdefinition}{def:p09-gpss-canonical-prefix}{The canonical positive-chain event test}
For a chain requiring point labels \(\Theta_1,\ldots,\Theta_m\)
and the preceding gap labels \(\Phi_0,\ldots,\Phi_{m-1}\),
with no final gap requirement, define future formulas
\[
 \Lambda_m=\top,\qquad
 \Lambda_i=\Phi_i\mathbin{\mathsf U_s}
                   (\Theta_{i+1}\land\Lambda_{i+1})
                       \quad(0\leq i<m).
\]
For bounds \(a<b\) within the root interval, define
\[
 \mu_1=\operatorname{First}_{\Theta_1\land\Lambda_1}(a,b),\qquad
 \mu_i=\operatorname{First}_{\Theta_i\land\Lambda_i}(\mu_{i-1},b)
                       \quad(2\leq i\leq m).
\]
The Boolean event test is
\[
 V_m(a,b)=
   \bigwedge_{i=1}^{m}\operatorname{Def}(\mu_i)
   \land\neg\operatorname{Def}
          (\operatorname{Last}_{\neg\Phi_0}(a,\mu_1)).
\]
For \(m=0\), set \(V_0=\mathsf{true}\), using no event
terms. Each use of \(V_m\) retains its own displayed point and
gap lists. The \(\Lambda_i\) formulas do not contain the
bounds \(a,b\) as parameters.
\end{hbdefinition}

\begin{hblemma}{thm:p09-gpss-canonical-prefix}{proved=annex:p9.gpss.canonical.prefix}{Correctness and minimality of the canonical chain}
On a bounded interval in a discrete complete order, \(V_m(a,b)\)
holds exactly when there is a strict chain of the required \(m\)
points and preceding gaps inside \((a,b)\). When it holds,
the selected points \(\mu_i\) satisfy those conditions and are
coordinatewise no later than the points of any other satisfying
chain. Future viability in \(\Lambda_i\) may inspect beyond
\(b\); the bounded event selectors still enforce every actual
witness to be strictly before \(b\).
\end{hblemma}

\begin{hbdefinition}{def:p09-gpss-determining-events}{The finite determining-list recursion}
Write a bounded decomposition as
\(D_W(a,b)\), where
\(W=(\Theta_1,\ldots,\Theta_r;\Phi_0,\ldots,\Phi_r)\).
Its rank is the number \(r\) of internal points. For \(0\leq j<r\),
let \(W_j\) be its prefix with \(j\) points and gap labels
\(\Phi_0,\ldots,\Phi_j\). Let \(V^{j}(a,b)\) be the
canonical positive-chain test for points \(\Theta_{j+1},\ldots,
\Theta_r\) with preceding gaps \(\Phi_j,\ldots,\Phi_{r-1}\);
thus \(V^{r}=\mathsf{true}\).

For rank zero define the decision expression
\[
 T_W(a,b)=\neg\operatorname{Def}
                    (\operatorname{Last}_{\neg\Phi_0}(a,b)).
\]
For positive rank put \(\ell=\operatorname{Last}_{\neg\Phi_r}(a,b)\)
and define
\[
\begin{aligned}
 T_W(a,b)={}&
  \bigl(\neg\operatorname{Def}(\ell)\land V^0(a,b)\bigr)\\
 &\lor\Bigl(\operatorname{Def}(\ell)\land
        \bigl[\ \bigvee_{j=0}^{r-1}
          \bigl(T_{W_j}(a,\ell)\land
             \operatorname{Def}(\operatorname{Test}_{\Phi_j}(\ell))
             \land V^{j}(\ell,b)\bigr)\\
 &\hspace{31mm}\lor\bigvee_{j=1}^{r}
          \bigl(T_{W_{j-1}}(a,\ell)\land
             \operatorname{Def}(\operatorname{Test}_{\Theta_j}(\ell))
             \land V^{j}(\ell,b)\bigr)\ \bigr]\Bigr).
\end{aligned}
\]
The first inner disjunction puts \(\ell\) in a gap; the second
puts it at an internal point. Form the event list \(E_W\) from
all terms in this expression and every subterm, including \(x,y\).
Each recursive decomposition has smaller rank. Identical endpoint
subterms can be shared as a finite acyclic expression graph.
\end{hbdefinition}

\begin{hblemma}{thm:p09-gpss-determining-events}{proved=annex:p9.gpss.determining.events}{The determining list is finite and decides the decomposition}
The recursion terminates and gives a finite event list whose definedness
statuses and values determine the truth of \(D_W(x,y)\). More
precisely, on actual endpoints \(x<y\),
\[
                             D_W(x,y)\longleftrightarrow T_W(x,y).
\]
All labels in the list are finite future-only \(L(U)\) formulas.
The construction includes the last-gap failure being absent, lying
in any gap, or coinciding with any internal witness. No label is
tested at either external endpoint.
\end{hblemma}

\begin{hbtheorem}{thm:p09-gpss-bounded-complement}{proved=annex:p9.gpss.bounded.complement}{Constructive relative complement on bounded intervals}
For every GPSS decomposition \(D_W(x,y)\), a finite procedure
constructs decompositions \(C_1(x,y),\ldots,C_N(x,y)\), with
the same permitted future-only label class, such that over every
discrete complete model
\[
             x<y\land\neg D_W(x,y)
                        \longleftrightarrow\bigvee_{i=1}^{N}C_i(x,y).
\]
The procedure enumerates all defined/undefined assignments to the
finite event list and all weak equality/order types of the defined
values, with \(x<y\). It compiles every consistent case into one
decomposition by the following exact constraints:

\begin{enumerate}
\item A defined first occurrence has its label true and has no earlier
      occurrence between its lower bound and its value; a defined last
      occurrence has no later occurrence before its upper bound.
\item An undefined first or last occurrence with defined ordered bounds
      has no occurrence anywhere between those bounds. Undefined or
      unordered bounds force an undefined result without this extra test.
\item A defined test equals its argument and has its label true; an
      undefined test with a defined argument has the label false there.
\item Universal interval constraints are transferred to every merged
      gap and every inserted point they cover. All point constraints
      are conjoined at their value's equality class.
\end{enumerate}

Retain exactly the cases where \(T_W\) is false. Impossible order
or status cases contribute no disjunct; contradictory label cases
may remain as unsatisfiable decompositions. Empty output and empty
open intervals are included. This supplies the finite determining
and back-translation steps of the event-expression outline for
bounded intervals
\cite[Proposition~2.4, p.~171]{GabbayPnueliShelahStavi1980}.
\end{hbtheorem}

\begin{hbtheorem}{thm:p09-gpss-suffix-complement}{proved=annex:p9.gpss.suffix.complement}{Complementation of future rays, including last points}
For an actual lower endpoint \(x\), the complement of every
\(D(x,+\infty)\) has a finite disjunction of exact GPSS
decompositions on \((x,+\infty)\). This includes all the
allowed global cofinal-occurrence conjuncts.

Take its already proved future expression \(B(x)\) from
\cref{thm:p09-gpss-suffix-translation} and put \(E=\neg B\).
Its outer Boolean construction has only constants and strict
future modalities as leaves. At a last point evaluate those leaves
by \(F A=\mathsf{false}\), \(G A=\mathsf{true}\), and
\(A\mathbin{\mathsf U_s}C=\mathsf{false}\), obtaining a
Boolean value \(e_{\mathrm{last}}\). For a point having a next point,
form the future label \(L_E\) by Boolean recursion using
\[
\begin{aligned}
 L_{F A}&=A\lor F A,\\
 L_{G A}&=A\land G A,\\
 L_{A\mathbin{\mathsf U_s}C}
       &=C\lor\bigl(A\land(A\mathbin{\mathsf U_s}C)\bigr).
\end{aligned}
\]
One output decomposition has one internal point labelled \(L_E\),
first gap labelled false and final gap labelled true. If
\(e_{\mathrm{last}}\) is true, add the zero-point decomposition
whose only gap is labelled false. These are all the output cases.
\end{hbtheorem}

\begin{hbtheorem}{thm:p09-gpss-prefix-complement-obstruction}{proved=annex:p9.gpss.prefix.complement.obstruction}{Why formal left endpoints require a different claim}
The same future-labelled prefix class is not closed under
complement, even over \(\langle\mathbb Z,<\rangle\).
The legal decomposition
\[
 D(-\infty,y)=\exists z<y\ \forall u<z\ q_0(u)
\]
has one internal point labelled true, initial gap labelled \(q_0\),
and final gap labelled true. Its complement has no finite disjunction
of prefix decompositions whose labels remain in \(L(U)\).
The analogous whole-order decomposition
\(D(-\infty,+\infty)=\exists z\,\forall u<z\,q_0(u)\)
has the same failure, even with the permitted future-cofinal conjuncts.

Thus the future proof uses bounded intervals and future rays with an
actual lower endpoint. A claim about all formal-endpoint configurations
in the unchanged future-label class would be false; changing to a
past-label or enlarged coinitial grammar would be a separate construction.
\end{hbtheorem}

\begin{hbtheorem}{thm:p09-gpss-discrete-future-completeness}{proved=annex:p9.gpss.discrete.future.completeness}{Completion of the discrete-complete future construction}
On discrete complete orders, every member of
\(\operatorname{Bet}_G(x,y)\) on \(x<y\), and every
member of \(\operatorname{Aft}_G(x)\), has a finite
disjunction of the exact GPSS decompositions. Here \(x\)
is an actual point variable. Consequently every source future
\(L_1\)-formula at \(x\), with quantifiers restricted to
\([x,+\infty)\), has an equivalent \(L(U)\)-formula at
every point of every discrete complete model. This is Theorem~2.2.
Every \(L_1\)-formula with only \(x\) free therefore also has
an equivalent \(L(U)\)-formula at point zero of every
\(\omega\)-model, the exact conclusion of Theorem~2.1.

The bounded failure on arbitrary complete orders is proved in
\cref{thm:p09-gpss-complete-bounded-obstruction}; the enlarged local
boundary tools and positive-chain construction are
\cref{thm:p09-gpss-infimum-case-compiler,thm:p09-gpss-complete-positive-chain}.
The discrete \(U\)-only conclusion above is not asserted beyond
its stated order class
\cite[Theorems~2.1--2.2, pp.~170--171]{GabbayPnueliShelahStavi1980}.
\end{hbtheorem}

\begin{hbexample}{def:p09-gpss-last-failure-example}{A last failure may be the witness itself}
On \(\{0,1,2,3\}\), let the endpoints be zero and three.
Take a rank-one word with first gap always true, final gap false
only at two, and its point label true only at two. Its last final-gap
failure is \(\ell=2\), and the decomposition is true with
its witness equal to \(\ell\). The final gap excludes that
witness, so the failed gap label there is harmless. If the point
label is instead true only at one, the decomposition is false.
Both outcomes are represented by the equality and gap cases of the
finite determining expression. On an adjacent endpoint interval,
every positive-rank decomposition is false and its relative complement
is true.
\end{hbexample}
